\documentclass[11pt]{article}
\usepackage[margin=1in]{geometry}
\usepackage{hyperref}
\usepackage{enumitem}
\title{Interactive Word Wall Build}
\author{Piter Garcia}
\date{\today}
\begin{document}
\maketitle

\section*{Grade and Class Match}
\textbf{Grade:} Grade 1\\
\textbf{Likely blocks:} Day 3 · 12:25-1:15 · Gargana; Day 3 · 2:15-3:05 · Polfleit; Day 5 · 12:25-1:15 · Montgomery

\section*{Lesson Focus}
Vocabulary, visual communication, and classroom making through an interactive build.

\section*{Learning Targets}
\begin{itemize}[leftmargin=*]
\item Students connect a word or concept to a physical or visual representation.
\item Students work with peers to assemble a simple shared product.
\item Students explain what their piece represents.
\end{itemize}

\section*{Materials}
\begin{itemize}[leftmargin=*]
\item Word cards
\item Maker materials
\item Display space
\end{itemize}

\section*{Suggested Lesson Flow}
\begin{enumerate}[leftmargin=*]
\item Open with the exact device, tool, or routine students need in order to start successfully.
\item Model one example task before students begin independent or partner work.
\item Give students time to build, code, design, or document their work while circulating for support.
\item Close with a short share-out, reflection, or debugging discussion.
\end{enumerate}

\section*{Source Notes}
This lesson packet is enriched with transcript-derived lesson notes, the school schedule roster, and official starting-point resources. See the paired Markdown guide for clickable links.

\bibliographystyle{plain}
\bibliography{interactive-word-wall-build}
\end{document}
